\section{The admissible cone}
\label{sec:hessian-cone}
\index{Hessian equation!admissible cone}\index{Garding cone@G\aa rding cone}

For $\lambda\in\mathbb R^n$, let $S_j(\lambda)$ be the elementary
symmetric polynomial of degree $j$, put $\sigma_j=S_j/\binom nj$ and
$\sigma_0=1$, and set
\[
 \Gamma_m=\{\lambda:S_j(\lambda)>0,\ 1\leq j\leq m\}.
\]
A real $(1,1)$-form belongs to $\Gamma_m(\omega)$ if its eigenvalues
relative to $\omega$ belong to $\Gamma_m$. Write $C_{n,m}$ for the
corresponding closed cone of Hermitian matrices. The algebraic assertions
below also apply when $m=n$.

We use the classical G\aa rding facts: the cone $C_{n,m}$ is convex,
contains the nonnegative matrices, and is preserved by adding them;
adding a positive matrix puts a matrix of $C_{n,m}$ in its interior.
The normalized root $F=\sigma_m^{1/m}$ is concave and increasing on
$\Gamma_m$. Polarized products of $k$ forms in $\Gamma_k(\omega)$ with
$\omega^{n-k}$ are positive, and on $\Gamma_k$, hence by continuity on $\overline{\Gamma_k}$, Maclaurin's inequalities
give $\sigma_j^{1/j}\geq\sigma_k^{1/k}$ for $1\leq j\leq k$. If $\alpha$
has eigenvalues $\lambda$ relative to $\omega$, then
$\alpha^j\wedge\omega^{n-j}=\sigma_j(\lambda)\,\omega^n$; this is the
normalization of all Hessian measures in the chapter.

A continuous real function $v$ on $X$ is \emph{admissible} if
\begin{equation}\label{eq:hessian-continuous-admissibility}
 (\omega+\ddc v)\wedge\gamma_1\wedge\cdots\wedge\gamma_{m-1}
                                  \wedge\omega^{n-m}\geq0
\end{equation}
in distributions for every smooth collection
$\gamma_1,\ldots,\gamma_{m-1}\in\Gamma_m(\omega)$.
For smooth $v$ this is equivalent, by G\aa rding duality, to
$\omega+\ddc v\in\overline{\Gamma_m(\omega)}$ at every point; a smooth
$v$ is \emph{strictly admissible} if $\omega+\ddc v\in\Gamma_m(\omega)$
everywhere. Maxima and convex combinations preserve continuous
admissibility. To check the maximum property, fix
$\gamma_1,\ldots,\gamma_{m-1}$ in
\eqref{eq:hessian-continuous-admissibility}. By polarized positivity,
the left-hand side is $(Lv+g)\,\omega^n$ for a linear elliptic operator
$L$ with smooth coefficients and a smooth function $g$. Subtracting a
local smooth solution of $Lh=-g$ reduces the assertion to the classical
fact that the maximum of two continuous distributional subsolutions of
$L$ is again a subsolution.

For a bounded upper semicontinuous function $v$, admissibility means that
every \emph{upper test} $q$ at a point $x$, that is, every $C^2$ function
$q$ near $x$ with $q\geq v$ near $x$ and $q(x)=v(x)$, satisfies
\begin{equation}\label{eq:hessian-bounded-upper-tests}
 \omega(x)+\ddc q(x)\in\overline{\Gamma_m(\omega(x))}.
\end{equation}
For continuous $v$, this agrees with
\eqref{eq:hessian-continuous-admissibility}: G\aa rding duality reduces
both conditions to subsolution inequalities for smooth strictly elliptic
linear operators, and we use the classical equivalence of viscosity and
distributional subsolutions for such operators. In the bounded class,
pairing \eqref{eq:hessian-bounded-upper-tests} with $\omega^{n-1}$
makes $v$, up to a local smooth function, a viscosity subsolution, and
hence a subharmonic function, for such an operator. As for ordinary
subharmonic functions, $v(x)$ is then the limit of the essential
suprema of $v$ on small coordinate balls centered at $x$. Thus $v$ is
the \emph{canonical representative} of its class in $L^1(X)$, and
this representative is determined by the distribution defined by $v$.
\index{Admissible potential!bounded upper tests}\index{Admissible potential!canonical representative}

The first estimate controls a scalar shift of the cone. It will give
lower bounds for every coefficient of the linearized Hessian operator.

\Needspace{14\baselineskip}
\begin{lemma}\label{hessian:hessian-cone-shift}
For $\lambda\in\Gamma_m$, put
\[
 e_m=\frac{\sigma_m(\lambda)}{m\sigma_{m-1}(\lambda)}.
\]
Then $\lambda-t\mathbf1\in\Gamma_m$ for $0\leq t<e_m$. If $m\geq2$, the conclusion also holds for $t=e_m$. For $m=1$, the endpoint lies on the boundary.
\end{lemma}
\begin{proof}
The polynomial
\[
 P(x)=\sigma_m(\lambda+x\mathbf1)
     =\sum_{j=0}^m\binom mj\sigma_j(\lambda)x^{m-j}
\]
is a positive multiple of the $(n-m)$th derivative of $\prod_i(x+\lambda_i)$. Hence all its roots are real, by Rolle's theorem. Its coefficients are positive, so all its roots are strictly negative. Write $P(x)=\prod_{i=1}^m(x+r_i)$ with $r_i>0$. Logarithmic differentiation at zero gives
\[
 \sum_{i=1}^m r_i^{-1}
       =\frac{m\sigma_{m-1}(\lambda)}{\sigma_m(\lambda)}
       =e_m^{-1}.
\]
Thus $\min_i r_i\geq e_m$, with strict inequality if $m\geq2$. For $x> -\min_i r_i$, every derivative of $P$ is positive: differentiating its product expansion produces sums of products of positive factors. Since
\[
 P^{(m-j)}(x)=\frac{m!}{j!}\sigma_j(\lambda+x\mathbf1),
 \qquad 1\leq j\leq m,
\]
all the required symmetric polynomials are positive at $x=-t$. When $m=1$, the sole root is $-e_1$.
\end{proof}

\Needspace{10\baselineskip}
\begin{lemma}\label{hessian:hessian-mixed-trace}
If $\alpha,\gamma\in\Gamma_m(\omega)$, then
\[
 \frac{\alpha^{m-1}\wedge\gamma\wedge\omega^{n-m}}{\omega^n}
 \frac{\alpha\wedge\omega^{n-1}}{\omega^n}
 \geq\frac1{m!}
 \frac{\alpha^m\wedge\omega^{n-m}}{\omega^n}
 \frac{\gamma\wedge\omega^{n-1}}{\omega^n}.
\]
The assertion extends to $\gamma\in\overline{\Gamma_m(\omega)}$ by continuity.
\end{lemma}
\begin{proof}
Write $\sigma_k=\alpha^k\wedge\omega^{n-k}/\omega^n$, and
$M_k=\alpha^{k-1}\wedge\gamma\wedge\omega^{n-k}/\omega^n$.
Let $1\leq k<m$. For $0\leq t<\sigma_{k+1}/((k+1)\sigma_k)$, \cref{hessian:hessian-cone-shift}, with $k+1$ in place of $m$, puts $\alpha-t\omega$ in $\Gamma_{k+1}(\omega)$. G\aa rding positivity for the $k+1$ forms $\alpha-t\omega,\alpha,\ldots,\alpha,\gamma$ of $\Gamma_{k+1}(\omega)$ gives $(\alpha-t\omega)\wedge\alpha^{k-1}\wedge\gamma\wedge\omega^{n-k-1}\geq0$, that is,
\[
 M_{k+1}\geq tM_k
 \quad\left(0\leq t<\frac{\sigma_{k+1}}{(k+1)\sigma_k}\right).
\]
Let $t$ increase to the displayed endpoint. Iterating, starting with $M_1=\gamma\wedge\omega^{n-1}/\omega^n$, and using $\prod_{k=1}^{m-1}\sigma_{k+1}/((k+1)\sigma_k)=\sigma_m/(m!\sigma_1)$, yields
\[
 M_m\geq\frac{\sigma_m}{m!\sigma_1}M_1.
\]
All denominators are positive. Approximating a form in the closed cone by $\gamma+\delta\omega$, $\delta\downarrow0$, gives the last assertion.
\end{proof}

At a diagonal Hermitian matrix $\Lambda$ in the interior of $C_{n,m}$, order its
eigenvalues as $\lambda_1\geq\cdots\geq\lambda_n$ and write
\[
 f=F(\Lambda),\quad F_i=\partial_{\lambda_i}F,\quad
 S=\operatorname{tr}\Lambda,\quad \mathcal F=\sum_iF_i,
 \quad Q=\sum_iF_i\lambda_i^2.
\]
The next inequalities use these derivatives of the normalized root,
rather than derivatives of the polynomial $S_m$.

\Needspace{10\baselineskip}
\begin{lemma}\label{hessian:weighted-cone-trace}
If $M\in C_{n,m}$, then
\[
 W_M:=\sum_iF_iM_{i\bar i}\geq\frac{f}{m!S}\operatorname{tr}M,
 \qquad F_i\geq\frac{f}{m!S},\qquad
 |M_{1\bar1}|\leq\operatorname{tr}M.
\]
\end{lemma}
\begin{proof}
In the notation of the proof of \cref{hessian:hessian-mixed-trace}, with $\alpha$ and $\gamma$ the forms with matrices $\Lambda$ and $M$, differentiation gives $DF_\Lambda(M)=fM_m/\sigma_m$. The estimate $M_m\geq\sigma_m M_1/(m!\sigma_1)$, with $M_1=\operatorname{tr}M/n$ and $\sigma_1=S/n$, proves the first inequality. The positive semidefinite rank-one coordinate projections belong to $C_{n,m}$; taking $M$ to be each projection proves the second inequality.
Since $m\geq2$, $C_{n,m}$ is contained in $C_{n,2}$, so the eigenvalues $\mu$ of $M$ satisfy
\[
 \operatorname{tr}M\geq0,\qquad
 \sum_i\mu_i^2=(\operatorname{tr}M)^2-2S_2(\mu)
                   \leq(\operatorname{tr}M)^2.
\]
Thus the operator norm of $M$ is at most its trace, which proves the last inequality. The wedge pairing computes the derivative in the diagonal frame of $\Lambda$, regardless of the eigenvectors of $M$.
\end{proof}

\begin{lemma}\label{hessian:quadratic-coercivity}
If $0<f_0\leq f\leq f_1$, then
\[
 Q:=\sum_iF_i\lambda_i^2
       \geq\frac{f_0}{m}S-\frac{n-m}{m}f_1^2,
 \qquad \mathcal F\geq1.
\]
\end{lemma}
\begin{proof}
Write $(\lambda|i)$ for $\lambda$ with its $i$th coordinate omitted. The derivative of the normalized root is $F_i=fS_{m-1}(\lambda|i)/(mS_m)$. The splitting identity $S_k(\lambda)=\lambda_iS_{k-1}(\lambda|i)+S_k(\lambda|i)$, applied with $k=m$ and $k=m+1$, gives
\[
 Q=\frac f m S-\frac{(m+1)f}{mS_m}S_{m+1}.
\]
For $m=n$, set $S_{n+1}=0$. If $S_{m+1}\leq0$, the second term is nonnegative. Otherwise $\lambda\in\Gamma_{m+1}$ and Maclaurin gives $S_{m+1}\leq\binom n{m+1}f^{m+1}$. Since $S_m=\binom nm f^m$, the subtracted term is at most $(n-m)f^2/m$. This proves the first assertion. Summing the derivatives gives $\mathcal F=f\sigma_{m-1}/\sigma_m\geq1$, again by Maclaurin.
\end{proof}

When $m<n$, admissible forms need not be semipositive. Positivity of a
mixed Hessian measure therefore does not directly bound every coefficient
of its factors. For degrees below $m$, deleted-coordinate positivity
supplies the required absolute bound: if $S_{j;I}(\lambda)=S_j(\lambda|I)$
denotes $S_j$ with the coordinates indexed by $I$ omitted, then
$S_{j;I}(\lambda)>0$ for $\lambda\in\Gamma_m$ and $j+|I|\leq m$, by
polarized G\aa rding positivity. By continuity, $S_{j;I}\geq0$ on
$\overline{\Gamma_m}$ in the same range.

\begin{lemma}\label{hessian:absolute-mixed-forms}
If $0\leq r<m$, $\alpha_1,\ldots,\alpha_r\in\Gamma_m(\omega)$,
and $\Xi$ is a smooth form of bidegree $(n-r,n-r)$, then
\[
 |\alpha_1\wedge\cdots\wedge\alpha_r\wedge\Xi|
 \leq C_{n,m,\Xi}(\alpha_1+\cdots+\alpha_r)^r\wedge\omega^{n-r}.
\]
For $r=0$ the right-hand side means $C_{n,m,\Xi}\omega^n$. The same
bound holds for $\alpha_1,\ldots,\alpha_r\in\overline{\Gamma_m(\omega)}$.
\end{lemma}
\begin{proof}
First take a single form $\alpha$, diagonal in an $\omega$-unitary frame,
with ordered eigenvalues $\lambda\in\Gamma_m$.
For $1\leq r<m$, repeated splitting and deleted-coordinate positivity
give $S_r(\lambda)\geq\lambda_1\cdots\lambda_r$.
The sum of the last $n-r$ eigenvalues is $S_{1;I}(\lambda)>0$, with
$I=\{1,\ldots,r\}$, so each negative
eigenvalue has modulus at most $(n-r)\lambda_{r+1}$.
Thus every product of $r$ distinct eigenvalues has modulus at most
\[
 (n-r)^r\lambda_1\cdots\lambda_r\leq(n-r)^r S_r(\lambda).
\]
Indeed, replace its negative factors by this bound and its positive
factors by the corresponding largest eigenvalues. In the diagonal frame
only complementary diagonal coefficients of $\Xi$ contribute to
$\alpha^r\wedge\Xi$. This proves
$|\alpha^r\wedge\Xi|\leq C_{n,m,\Xi}\alpha^r\wedge\omega^{n-r}$.

Apply this bound to the polynomial
$(\alpha_1+t_2\alpha_2+\cdots+t_r\alpha_r)^r\wedge\Xi$,
with $0\leq t_i\leq1$. Since
$\beta=\alpha_1+t_2\alpha_2+\cdots+t_r\alpha_r$ lies in
$\Gamma_m(\omega)$, its absolute value is at most
$C_{n,m,\Xi}\beta^r\wedge\omega^{n-r}$. Expanding
$(\alpha_1+\cdots+\alpha_r)^r=(\beta+\sum_{i\geq2}(1-t_i)\alpha_i)^r$
and using polarized positivity bounds this by the stated right-hand side.
A fixed invertible interpolation matrix on this finite-dimensional
polynomial space bounds each coefficient by its supremum on the cube.
The coefficient of $t_2\cdots t_r$ is a fixed positive multiple of
$\alpha_1\wedge\cdots\wedge\alpha_r\wedge\Xi$.
This proves the mixed bound. The closed-cone case follows by adding
$\varepsilon\omega$ to each form and taking the limit.
\end{proof}

The next coefficient bound also controls a torsion form whose
diagonal coefficients can have either sign. It compares each
coefficient $S_{q;ij}$ with two deleted coordinates with the
nonnegative coefficients $S_{q;i}$ and $S_{q;j}$ of the associated
gradient pairing.

\Needspace{14\baselineskip}
\begin{lemma}\label{hessian:deleted-newton}\index{Newton inequality!deleted coordinates}
If $\lambda\in\overline{\Gamma_m}$, $1\leq q\leq\min\{m-1,n-2\}$, and $i\ne j$,
then, with $N=n-2$ and
\[
 c_{N,q}=\frac{q(N-q)}{(q+1)(N-q+1)}<1,
\]
one has
\[
 |S_{q;ij}(\lambda)|^2
 \leq\frac{S_{q;i}(\lambda)S_{q;j}(\lambda)}{1-c_{N,q}}.
\]
\end{lemma}
\begin{proof}
First prove the unnormalized Newton inequality for an arbitrary real
list $\mu$ of length $N$. For $1\leq q<N$, put
$R(x)=\prod_{a=1}^N(x+\mu_a)$ and
$G=R^{(N-q-1)}$. Repeated Rolle's theorem makes $G$ a real-rooted
polynomial of degree $d=q+1$. Away from its roots $r_a$, logarithmic
differentiation and Cauchy--Schwarz give
\[
 G'^2-GG''=G^2\sum_{a=1}^d\frac1{(x-r_a)^2}
 \geq \frac1dG^2\left(\sum_{a=1}^d\frac1{x-r_a}\right)^2
 =\frac1dG'^2.
\]
Continuity extends this polynomial inequality to every real $x$.
At zero,
\[
 G=(N-q-1)!S_{q+1}(\mu),\quad
 G'=(N-q)!S_q(\mu),\quad
 G''=(N-q+1)!S_{q-1}(\mu).
\]
Substituting these values in $(1-1/d)G'^2\geq GG''$ gives
$S_{q-1}(\mu)S_{q+1}(\mu)\leq c_{N,q}S_q(\mu)^2$.
No sign restriction on $\mu$ or on these coefficients was used.
For $q=N$, $S_{N+1}(\mu)=0$, so the same inequality holds
with $c_{N,N}=0$.

Apply this to the list with $i,j$ deleted. Put
$A=S_{q-1;ij}$, $B=S_{q;ij}$, $C=S_{q+1;ij}$ and
$D_i=S_{q;i}$, $D_j=S_{q;j}$. Deleted-coordinate positivity gives
$A,D_i,D_j\geq0$: the required degrees are $q+1\leq m$.
Splitting in the two omitted coordinates gives the exact identity
\[
 D_iD_j=A S_{q+1}(\lambda)+B^2-AC.
\]
Since $S_{q+1}(\lambda)\geq0$ and $AC\leq c_{N,q}B^2$, this is
at least $(1-c_{N,q})B^2$.
\end{proof}

A smooth maximum preserves the cone condition and can be chosen
monotone in its smoothing scale as well as in its two arguments.

\Needspace{16\baselineskip}
\begin{lemma}\label{hessian:regularized-maximum}\index{Regularization!Hessian potentials}
If smooth uniformly bounded admissible $u_j\downarrow u$ and
$v_j\downarrow v$ are given, there are smooth uniformly bounded
admissible $w_j\downarrow\max\{u,v\}$. For any chosen
$\epsilon_j>0$ decreasing to zero, they can be chosen to satisfy
$w_j=u_j$ on $\{u_j-v_j>\epsilon_j\}$ and
$w_j=v_j$ on $\{v_j-u_j>\epsilon_j\}$.
The pointwise maximum of two bounded admissible functions is bounded
admissible. For smooth functions $f,g$ on an open subset of $X$ with
$\ddc f,\ddc g\in\overline{\Gamma_m(\omega)}$, the regularized
maximum $M_\epsilon(f,g)$ constructed below satisfies the same
unshifted cone condition.
\end{lemma}
\begin{proof}
Choose a nonnegative even smooth probability density $\zeta$
supported in $[-1,1]$, and set
\[
 r_\epsilon(t)=\int |t-\epsilon s|\zeta(s)\,\mathrm ds,
 \qquad M_\epsilon(x,y)=\tfrac12(x+y+r_\epsilon(x-y)).
\]
The function $r_\epsilon$ is smooth and convex,
$|r_\epsilon'|\leq1$, and $r_\epsilon(t)=|t|$ for
$|t|\geq\epsilon$. Jensen's inequality and the triangle inequality
give
\[
 \max\{x,y\}\leq M_\epsilon(x,y)
                     \leq\max\{x,y\}+\epsilon/2.
\]
Its partial derivatives in $x,y$ are nonnegative. For fixed $t$,
$r_\epsilon(t)$ is an even convex function of $\epsilon$, since
$\zeta$ is even. It is therefore nondecreasing for $\epsilon\geq0$.
Consequently $w_j=M_{\epsilon_j}(u_j,v_j)$ decreases pointwise to
$\max\{u,v\}$ and has the asserted equalities and bounds.

Put $h_j=u_j-v_j$ and
$a_j=(1+r_{\epsilon_j}'(h_j))/2$. The chain rule gives
\[
 \omega+\ddc w_j
 =a_j(\omega+\ddc u_j)+(1-a_j)(\omega+\ddc v_j)
       +\tfrac12r_{\epsilon_j}''(h_j)
                    \mathrm dh_j\wedge\dc h_j.
\]
Here $0\leq a_j\leq1$, and the last form is semipositive.
Convexity of the cone and its stability under semipositive addition
prove admissibility. For smooth local functions satisfying the unshifted
cone condition, omit $\omega$ from this chain rule; the same argument
proves preservation of that condition.

Finally, let $u,v$ be bounded admissible. Their pointwise maximum is
bounded and upper semicontinuous, and an upper test for it at a point
is an upper test for whichever of $u,v$ attains the maximum there.
Hence $\max\{u,v\}$ satisfies \eqref{eq:hessian-bounded-upper-tests};
in particular, it is its own canonical representative.
\end{proof}

\Needspace{11\baselineskip}
\begin{lemma}\label{hessian:canonical-order}\index{Admissible potential!pointwise order}
If bounded admissible $u,v$ satisfy $u\leq v$ almost everywhere
with respect to $\omega^n$, then $u\leq v$ everywhere.
\end{lemma}
\begin{proof}
By \cref{hessian:regularized-maximum}, $w=\max\{u,v\}$ is bounded
admissible. Since $w=v$ almost everywhere, $w$ and $v$ define the same
distribution. As bounded admissible functions, both are canonical
representatives, and a canonical representative is determined by its
distribution. Thus $w=v$ everywhere, that is, $u\leq v$.
\end{proof}
